\section{Supporting Workflows}
\label{sec:supporting}

In addition to the four primary workflows, the framework includes two auxiliary workflows for dataset management and synthetic data generation.

\subsection{Workflow 6: Synthetic Data Generation}

This workflow acts as an \textbf{Automated Pipeline Validator} and data augmentation source. While modern Generative AI (like Stable Diffusion) is often used for data expansion, this module implements deterministic Digital Signal Processing (DSP) and statistical simulation to serve two critical architectural roles:

\begin{enumerate}
    \item \textbf{Pipeline Sanity Check}: Before fine-tuning on large external datasets (which may contain corruption or domain shifts), the system generates "perfect" synthetic samples (e.g., clean sine waves, white noise) to validate the integrity of the training loop and the model's capacity to learn basic patterns.
    \item \textbf{Sensor Simulation}: For embedded sensors (accelerometers, microphones), mathematical models (chirps, drift simulation) often provide more controllable test cases than scarce real-world anomaly data.
\end{enumerate}

Key features include:
\begin{itemize}
    \item \textbf{Requirement Collection}: Extracts signal parameters (frequency, duration, noise floor) from user specifications via LLM parsing.
    \item \textbf{DSP Generation Engine}: Uses NumPy-based vectorized operations to generate time-series data with precise control over signal-to-noise ratio (SNR).
    \item \textbf{Standardized Output}: Produces artifacts in the exact \texttt{.npy} format expected by the Fine-Tuning Workflow (Workflow 5), confirming interface compatibility.
\end{itemize}

\subsection{Workflow 7: Dataset Selection and Format Conversion}

This workflow curates and converts publicly available datasets for embedded AI tasks:

\begin{itemize}
    \item \textbf{Dynamic Dataset Registry}: Queries a multi-tier catalog managed through external JSON files (\texttt{predefined\_datasets.json} and \texttt{dataset\_mapping.json}). This architecture decouples the knowledge base from the logic, enabling runtime updates of available datasets and their task-based recommendations.
    \item \textbf{Download and Cache}: Fetches datasets from Kaggle, UCI ML Repository, or GitHub
    \item \textbf{Format Standardization}: Converts to NumPy arrays with shape \texttt{(N, H, W, C)} for images or \texttt{(N, T, F)} for time series.
    \item \textbf{Generic Vision Processor}: Automatically discovers image files within heterogeneous directory structures, inferring labels from parent subdirectories. This feature enables zero-configuration processing of arbitrary vision datasets (e.g., Fruit-360) by handling archive extraction and automatic conversion to memory-efficient \texttt{.npy} formats with a capped sampling limit (up to 5,000 images) to prevent RAM exhaustion.
    \item \textbf{Train/Val/Test Splitting}: Applies stratified splits (default 70/15/15).
    \item \textbf{Audio Preprocessing}: Converts audio datasets to spectrograms via STFT for vision model compatibility.
\end{itemize}

The converted dataset is saved to \texttt{./dataset\_output/} and path is recorded in \texttt{MasterState} for consumption by Workflow 5.
